% FrontBackMatter/Abstract.tex
\addchaptertocentry{Хураангуй}

\vspace*{-1cm}
\vspace{0.8cm}

\begin{center}
{\headingfont\Large \ttitle\par}
\bigskip
{\headingfont\Large \deptname\par}
\bigskip

{\large \shortname\par}
{\headingfont\Large Хураангуй\par}

{\normalsize \par}
\end{center}
\vspace{1cm}


\bigskip

Монгол Улсын жижиг, дунд бизнесүүд хоорондоо бараа захиалах, хүргэх үйлээ гол төлөв утас, мессеж, цаасаар гүйцэтгэсээр байна. Ийм гараар зохицуулдаг хэлбэр нь цаг алдагдуулж, алдаа үүсгэж, явцыг хянахад төвөгтэй болгож байна.

Төсөлд зуучлагч оператор шаардахгүй, олон байгууллага зэрэг үйлчлүүлдэг, хүргэгчид өөрсдөө даалгавраа сонгодог B2B хүргэлтийн нэгдсэн платформыг боловсруулсан. Систем нь хүргэгчдэд зориулсан гар утасны аппликейшн, борлуулагчдад зориулсан веб аппликейшн, үүлэн дээр ажилладаг backend гэсэн гурван үндсэн хэсгээс бүрдэнэ.

Системийн найдвартай байдал, аюулгүй байдал, хурдны үзүүлэлтийг туршилтаар үнэлсэн. Нэг даалгаврыг олон хүргэгч зэрэг авахаар өрсөлдөхөд давхардал үүсэхгүй, байгууллагуудын мэдээлэл бие биендээ алдагдахгүй, хүргэлтийг баталгаажуулах код нь аюулгүй ажиллаж байгаа нь нотлогдсон.

Судалгааны үр дүнд тус систем нь Монголын жижиг, дунд бизнесийн хүргэлтийн гол бэрхшээлүүдийг бодитоор шийдэж чадах нь батлагдлаа.


\vfill
